
\paragraph{a) The reflected direction.}

The derivation follows the geometric explanation in
\href{https://www.pbr-book.org/3ed-2018/Reflection_Models/Specular_Reflection_and_Transmission#SpecularReflection}{\textit{Physically Based Rendering}}.

The reflected direction can be computed using simple vector geometry.
The incident direction $\omega_i$, the surface normal $\mathbf{n}$,
and the reflected direction $\omega_r$ all lie in the same plane.
For specular reflection, the reflected direction makes the same
angle with the normal as the incident direction. Therefore,
\[
    \omega_i \cdot \mathbf{n}
    = \omega_r \cdot \mathbf{n}.
\]

We can decompose $\omega_i$ into two components: one parallel
to the surface normal and one perpendicular to it,
\[
    \omega_i = \omega_{\parallel} + \omega_{\perp}.
\]
Since $\mathbf{n}$ is a unit vector,
\[
\begin{aligned}
    \omega_{\parallel}
        &= (\omega_i \cdot \mathbf{n})\mathbf{n},\\
    \omega_{\perp}
        &= \omega_i - (\omega_i \cdot \mathbf{n})\mathbf{n}.
\end{aligned}
\]

For perfect reflection, the normal component stays the same,
while the perpendicular component changes sign. Thus,
\[
\begin{aligned}
    \omega_r
        &= \omega_{\parallel} - \omega_{\perp}\\
        &= \omega_{\parallel} - (\omega_i-\omega_{\parallel})\\
        &= -\omega_i + 2(\omega_i\cdot\mathbf{n})\mathbf{n}.
\end{aligned}
\]

\paragraph{b) The local direction.}

In local coordinates, the surface normal is
$\mathbf{n}=(0,0,1)^\top$. Substituting this into the
reflection equation gives
\[
\begin{aligned}
    \omega_r
        &= -\omega_i + 2(\omega_i\cdot\mathbf{n})\mathbf{n}\\
        &= -(\omega_x,\omega_y,\omega_z)^\top
           + 2\omega_z(0,0,1)^\top\\
        &= (-\omega_x,-\omega_y,\omega_z)^\top.
\end{aligned}
\]


\paragraph{c) The refracted direction.}

Since incident rays can arrive from either side of the surface,
we first determine the incident and transmitted refractive indices:
\[
(\eta_i,\eta_t)=
\begin{cases}
(\eta_{\mathrm{ext}},\eta_{\mathrm{int}}), & \omega_z\geq 0,\\
(\eta_{\mathrm{int}},\eta_{\mathrm{ext}}), & \omega_z<0.
\end{cases}
\]

Let $\eta=\eta_i/\eta_t$. From Snell's law,
\[
    \eta_i\sin\theta_i=\eta_t\sin\theta_t.
\]
Like in reflection all vectors lie in the same plane.
For unit directions, the sine of the angle with the normal
equals the length of the perpendicular component. Therefore,
\[
    \omega_{t,x}=-\eta\omega_x,\qquad
    \omega_{t,y}=-\eta\omega_y.
\]

Since $\omega_t$ must also be a unit vector,
\[
    \omega_{t,z}^2=1-\eta^2(1-\omega_z^2).
\]
If we assume that the refracted direction points into the opposite medium,
its normal component has the opposite sign to $\omega_z$:
\[
\omega_t=
\begin{pmatrix}
-\eta\omega_x\\
-\eta\omega_y\\
-\operatorname{sgn}(\omega_z)
\sqrt{1-\eta^2(1-\omega_z^2)}
\end{pmatrix}.
\]
